% start document content
\chapter{Preface}
\section{Project Description}
	Designing aerospace gearboxes traditionally requires long hand calculations, trial-and-error assembly sketches, and 
	several rounds of costly physical tests. The Gearbox Artificial Intelligence Tool (GAIT) project was created to show 
	and prove how an automated workflow can replace most of that effort while still meeting the stringent 
	weight and safety of rotorcraft transmissions. 
	
	A compact Python routine takes 31 engineering inputs—power, shaft speeds, material properties, and 
	gear-geometry choices—and computes the initial 
	dimensions of a two-stage helicopter gearbox (spiral-bevel input and planetary output).  

	The resulting designs are then evaluated in the Romax Concept (Student Edition) software, which validates the 
	weight and size of the Python output result. It delivers accurate mass, safety-factor, and kinematic data for 
	each configuration. From these simulation results we build a fast “prediction engine” that, after a 
	brief training phase, can estimate the total gearbox weight and size for any new set of design parameters 
	in seconds. The prediction engine matches the full Romax Concept (Student Edition) results to within $2\%$ of weight 
	(e.g., 364 kg predicted vs. 357.8 kg simulated) while cutting the evaluation time from hours to a few seconds.

\section{Project Value}
	The value of the Gearbox Artificial Intelligence Tool (GAIT) is to harnesses AI to automate the manual design generation and 
	testing processes, creating an intelligent and repeatable workflow that accelerates 
	design development and optimization system and frees engineers to focus on 
	high-value engineering work.
		\begin{figure}[h]
			\centering
			\includegraphics[width=0.9\textwidth]{gearbox/gearbox1}
			\caption{GAIT}
			\label{fig:gearbox1}
		\end{figure}
		\FloatBarrier
		
\section{Internship Scope}
	The GAIT project's scope is the following: 
		\begin{itemize}
			\item Perform research into the engineering fundamentals and 
			real-world applications of gearbox layout focusing on design 
			and functionality. 
			\item If feasible, evaluate and develop a P-O-C on gearbox 
			parameter optimization using an AI-powered system on an existing 
			gearbox layout. 
			\item Explore the batch processing capabilities of Romax software
			 and evaluate the feasibility of automating gearbox layout 
			 validation workflows. 
		\end{itemize}

	\subsection{Internship Deliverables}
		The first phase of the project is the Summer 2026 cohort, with the following deliverables: 
		\begin{itemize}
			\item Conduct research and create a detailed report based on the key takeaways.
			\item Develop an AI system as a Proof of concept: \textit{With initial design parameters from provided resources to estimating optimal weight and size of the gearbox layout.}
			\item Explore ROMAX and its feasibility for design validation workflow.   
		\end{itemize}

		\subsection{Current Project State}
		The project has successfully established an analytical code parameterized by 31 key design variables, yielding to 
		give the corresponding weight and size of the gearbox design. In this project, 30,000-point dataset was generated 
		based on that code output and using quasi-random sampling to uniformly cover the N-dimensional design space. 

		Using the generated dataset, HistGradientBoosting Cascade Model was trained successfully to predicts gearbox weight 
		and volume with $R^2$ score of 0.987, bypassing standard gear snapping discontinuities.

		


